\documentclass[12pt,a4paper,oneside]{article}
\usepackage[brazil]{babel}
\usepackage[utf8]{inputenc}
\usepackage{olymp}

\setcounter{problem}{0}

\begin{document}

\begin{problem}{Relógio de atleta}{relogio}{2}

Uma empresa está desenvolvendo um novo relógio eletrônico para atletas de alto desempenho. Uma das funcionalidades do novo relógio é que ele vai medir a \textit{frequência cardíaca atual} (batidas do coração por minuto) e a \textit{capacidade de oxigenação atual} do atleta. O relógio também vai calcular e armazenar a \textit{frequência cardíaca em repouso} do atleta.

Os projetistas querem que o relógio emita um aviso para o atleta de que ele ou ela deve:\\[-18pt]

\begin{itemize}
\item \textit{diminuir} o ritmo do exercício se a frequência cardíaca atual é maior do que três vezes a frequência cardíaca em repouso ou a capacidade de oxigenação atual é menor do que 95;\\[-18pt]
\item \textit{aumentar} o ritmo do exercício se a frequência cardíaca atual é menor do que duas vezes frequência cardíaca em repouso e a capacidade de oxigenação atual é maior do que 97;\\[-18pt]
\item \textit{manter} o ritmo de exercício se nenhuma das condições anteriores ocorrer.\\[-18pt]
\end{itemize}

Dadas a frequência cardíaca em repouso, a frequência cardíaca atual e a capacidade de oxigenação atual obtidas pelo relógio, escreva um programa que produza a sugestão que o relógio deve emitir.

%\Notes
%
%Observações, se tiver.

\InputFile

A entrada é composta por três linhas, cada uma contendo um único número inteiro. A primeira linha contém o inteiro $R$, a frequência cardíaca em repouso do atleta. A segunda linha contém o inteiro $F$, a frequência cardíaca atual do atleta. A terceira linha contém o inteiro $C$, a capacidade de oxigenação atual do atleta. Restrições:  $35 \leq R \leq 100;\; 35 \leq F \leq 200;\; 50 \leq C \leq 100$.


\OutputFile

Seu programa deve produzir uma única linha, contendo uma única palavra, em letras minúsculas, que deve ser `\texttt{aumentar}', `\texttt{diminuir}', ou `\texttt{manter}', de acordo com os critérios estabelecidos acima.

% \Notes
% Note que a mensagem escrita não tem acento e tem um ponto final.

\Example

\begin{example}
\exmp{
60
119
98
}{
aumentar
}%
\end{example}

\begin{example}
\exmp{
60
190
100
}{
diminuir
}%
\end{example}

\begin{example}
\exmp{
50
100
95
}{
manter
}%
\end{example}

\begin{example}
\exmp{
50
100
94
}{
diminuir
}%
\end{example}


%Comentários sobre o exemplo, se necessário.
\footnotesize \textit{Problema da OBI 2020 - Olimpíada Brasileira de Informática}

\end{problem}

\end{document}
